\documentclass[a4paper]{article}
\usepackage[left=3.18cm,right=3.18cm,top=2.54cm,bottom=2.54cm]{geometry}
\usepackage{CJK}
\usepackage{indentfirst}
\usepackage{amsmath}
\usepackage{graphicx}
\usepackage{booktabs}
%%%%%%%%
%%%%%%%%
\usepackage{listings}
\begin{CJK}{UTF8}{gbsn}
\usepackage[usenames,dvipsnames,svgnames,table]{xcolor} % 用于code
\lstset{
	basicstyle=\ttfamily,
	xleftmargin=2pc, % 整体布局
	xrightmargin=2pc,
	backgroundcolor=\color{Lavender}, % 背景色
	rulecolor=\color{Silver}, % 边框颜色
	frame=single,% 边框
	commentstyle=\ttfamily\color{green!40!black},
	texcl=true,
}
\lstdefinelanguage{SAC} {
	keywords={SAC>}, % SAC提示符
	otherkeywords={SAC>}, % SAC提示符
	sensitive=true,
	comment=[l][{\color{green!40!black}}]{\#},
}
\lstnewenvironment{SACCode}{
	\lstset{
		language={SAC}, % 语言
		keywordstyle=\color{blue!70}\bfseries, % 关键字颜色
	}
}{}
\lstdefinelanguage{bash} {
	keywords={bash>}, % SAC提示符
	otherkeywords={bash>}, % SAC提示符
	sensitive=true,
	comment=[l][{\color{green!40!black}}]{\#},
}
\lstnewenvironment{bashCode}{
	\lstset{
		language={bash}, % 语言
		keywordstyle=\color{blue}, % 关键字颜色
	}
}{}


\linespread{1.3}
\addtolength{\parskip}{3pt}
\title{数据预处理{---从Ref\mbox{}Tek130数采到SAC}}
\author{汪晟}
\begin{document}
\begin{titlepage}
\begin{center}
\begin{figure}[ht!]
\centering
\vspace{4cm}
\includegraphics[width=0.25\textwidth]{./logo.JPG}
\end{figure}
%\textsc{\LARGE 中科院地质与地球物理研究所}\\[1.5cm]
\hrule
\vspace{0.4cm}
{\Huge \bfseries 数据预处理     }\\[0.5cm]
{\Large \bfseries {从Ref\mbox{}Tek130数采到SAC}}\\[1cm]
\hrule
\vspace{1cm}
\textsc{\large  作者：汪晟}\\[0.5cm]
\textsc{\large  版本：0.01}

\vfill
{\large \today}
\end{center}
\end{titlepage}

\tableofcontents

\clearpage\section{Ref格式数据的转换}
野外操作中，自130数采拷贝的原始数据是Ref格式，其单文件包含了多分量的数字记录以及部分台站与时间信息，可以使用PQL查看。此外，RefTek130数采亦生成log文件，实时记录了台站的状态信息，包括有：温度、GPS卫星数、电压等，可以使用newlogpeek.py查看。这些不同功能的软件均由PASSCAL（）提供。\par
根据需求，可以将Ref格式文件转换成其他格式的文件，相关程序均由PASSCAL提供。本文档依托项目为腾冲火山地区接受函数研究，需要数据转换成SAC格式。由Ref到SAC，需要借助中间文件格式，一般有两种方法较为常用：\par
\begin{center}
Ref $\to$ miniSeed $\to$ SAC\par
Ref $\to$ Segy     $\to$ SAC\par
\end{center}

本文档使用第二种方法，其命令行操作如下：\par
\begin{bashCode}
bash> ref2segy -f in.ref -l cas.file
bash> segy2sac in.ref					
\end{bashCode}

其中，ref2segy将Ref格式转换为Segy格式。-f参数声明Ref文件名或目录。-l参数声明130数采部分配置信息。一般情况下，原始Ref文件携有全部的数采配置信息，但是BUG时有发生，此时-l参数开启就很有必要。关于-l的详细介绍以及对应的130数采文件请参看ref2segy的manaul，在本文档的后续章节中，亦有提及，可供参考。\par
转换得到的SAC格式文件头信息缺失，这里引出下一节。

\clearpage\section{SAC格式头段信息简介}
除了时间序列记录之外，地震数据中的台站、事件等信息十分重要。本节承接上一节文件格式转换步骤，将简介SAC格式文件中的头段信息变量。\par
SAC头段信息变量包含了所有的数据参数信息，罗列接分类如下：\par
1. 基本变量 nvhdr nzyear nzjday nzhour nzmin nzsec nzmsec iztype iftype idep\par
2. 数据相关变量 npts delta odelta b e leven depmin depmax depmen scale xminimum xmaxmum yminimum ymaxmum nxsize nysize iqual isynth\par
3. 事件相关变量 kevnm evla evlo evel evdp ievreg ievtyp mag imagsrc imagtyp gcarc dist az bz o ko khole\par
4. 台站相关变量 kstnm knetwk istreg stla stlo stel stdp cmpaz cmpinc kcmpnm kstcmp lpspol\par
5. 震相相关变量 a ka f kf tn kn\par
6. 仪器想关变量 kinst iinst respn\par
7. 其他变量 usern kusern lovrok lcalda kdatrd\par

以上罗列中，×××变量为转换生产的SAC文件携有的有效头段信息。亦存在其他变量有值，但非有效，所以需要补全。\par
第一步需要补全的是台站相关变量。对任一数据记录，台站信息都是存在的且必须的。并非只有事件相关数据才可以用于研究，噪声数据的价值亦不可估量。\par
第二步需要补全的是事件相关变量。此操作只针对那些包含了时间记录的文件，所以需要先抽取出这些文件，再补全头信息。\par
第三步操作，需要补全部分震相等其他信息。\par
需要注意的是SAC头段中的时间类变量。由Ref转换得到的SAC文件携有的有效信息中有“参考时间”、“起始相对时间”、“结束相对时间”、“采样间隔”、“采样点数”。显然，这些时间信息倘若缺失，那么整个预处理都将失去意义。转换得到的SAC文件“参考时间”为数据首个采样点的绝对时间（伦敦时间），那么“起始相对时间”自然为0。在更为靠后的操作中，需要写入事件信息，并截断有效长度，此时“参考时间”会发生改变，相对类时间将会由SAC软件自动自动计算变化。总之，时间类变量的写入与更改，应当细心。\par
此外，SAC头段信息变量冗余，一些变量可由其他变量推导计算获得，这一特性减轻了补充头段信息变量的工作。\par

\clearpage\section{补全SAC文件头段台站信息变量}
上一节提及了对所有的文件记录，台站信息都是存在且必需的。本节将描述补全这一类信息。\par
SAC程序提供了更改任意头段信息变量的命令，如下：\par

\begin{SACCode}
SAC> r eg1.sacS
SAC> ch ARGV1 VAL1 ARGV2 VAL2 ...
SAC> w over
\end{SACCode}

第一行读SAC文件到内存中，第二行更改了内存中数据的头段变量值，第三行将内存中数据覆盖写入文件。本文档所有的SAC头段变量信息的赋值与更改均为以上操作。\par
本节中赋值更改的ARG范围列表如下：
\begin{table}[!htbp]
	\centering
	\caption{台站信息列表}\label{tab::StInfoTab}
	\begin{tabular}{llc}
	\toprule
	变量名	&解释				&	是否需要赋值\\\midrule
	kstnm	&台站名				&	是\\
	knetwk	&台网名 				&	是\\\midrule
	istreg	&台站地理区域\\
	stla	&台纬度				&	是\\
	stlo	&台经度				&	是\\
	stel	&台高度				&	是\\
	stdp	&台深度				&	是\\\midrule
	cmpaz	&通道方位角			&	是\\
	cmpinc	&通道倾角			&	是\\
	kcmpnm	&通道名				&	是\\
	kstcmp	&辅助名\\
	lpspol	&通道极性正负			&	是\\
	
	\bottomrule
	\end{tabular}
\end{table}
其中，通道名与通道角度之间的关系为：
\begin{table}[!htbp]
	\centering
	\caption{通道信息表}\label{tab::CmpInfo}
	\begin{tabular}{ccc}
	\toprule
	通道名	&cmpaz	&	cmpinc\\\midrule
	N		&0		&90\\
	E		&90		&90\\
	U		&0		&0\\
	\bottomrule
	\end{tabular}
\end{table}
\par
显然NEU为左手坐标系，lpspol变量为TRUE表明通道的正极性与NEU方向相同，为FALSE则相反。

\clearpage\section{补全SAC文件头段事件信息变量}
补全所有记录文件的台站信息之后，可挑选出包含有时间记录的文件。此时，当务之急是如何挑选出这类文件，并为这些文件补全事件类信息。\par
一种简单且可批量处理的方法是，给出台站与时间信息列表，组合成不同的“事件--台站”对，计算理论P波初至震相到时。由这些“事件--台站--到时”组合，挑选出包含了P波初至震相的记录文件。
更为复杂的情况是，需要提取出P波初至震相前后一段时长内的记录。那么这些时间记录片可能会跨不同文件，甚至缺失某些时段的记录文件。对此，可以事先给出P波初至震相前后时间偏移长度，计算出有效时间片的“起始绝对时间”和“终止绝对时间”，随后提取出所有与此时间区间存在交集的记录文件，合并即可。此步得到的“起始绝对时间”和“终止绝对时间”亦将用到后续的数据截断中。\par
获取了包含有事件记录的文件后，需要补全对应事件信息，以及理论计算的P波初至震相到时。变量列表如下：
\begin{table}[!htbp]
	\centering
	\caption{事件信息列表}\label{tab::EvInfoTab}
	\begin{tabular}{llc}
	\toprule
	变量名		&解释				&是否需要赋值\\\midrule
	kevnm   	&事件				&是\\\midrule
	evla		&事件纬度			&是\\
	evlo		&事件经度			&是\\
	evel		&事件高程			&是\\
	evdp		&事件深度			&是\\
	ievreg		&事件地理区域	\\\midrule
	ievtyp		&事件类型			&选\\\midrule
	mag			&震级				&是\\
	imagsrc		&震级信息来源			&是\\
	imagtyp		&震级类型			&是\\\midrule
	gcarc，dist	&大圆路径长			&当lcalda为TRUE时，将自动计算\\		
	az,bz		&前后方位角			&当lcalda为TRUE时，将自动计算\\\midrule
	t1			&理论P波初至相对时间	&是\\
	\bottomrule
	\end{tabular}
\end{table}

\clearpage\section{时间序列预处理}
补全SAC文件头段信息之后，接下来的预处理操作将针对时间序列展开。与头段信息不同的是，时间序列的更改不可逆，且不同的处理手段及先后顺序将直接影响到后续的数据研究分析。随意随性的操作将会带来假象，所以先于预处理，应有数据备份。\par
时间序列的预处理包括有：1.去毛刺、2去仪器相应、3.时间截断、4.坐标旋转、5.去均值与线性趋势、6.尖灭、7.滤波。简介：\par
%\begin{table}[!htbp]
%	\centering
%	\caption{预处理列表}\label{tab::ProcsTab}
%	\begin{tabular}{lll}
%	\toprule
%	编号	&预处理			&目的解释				\\\midrule
%	1	&去毛刺			&针对早期模拟记录中的尖峰和数据丢失，现代数字记录中已不存在此问题\\
%	2	&去仪器响应		&还原真实地动位移，并有带通滤波\\
%	3	&时间截断\\
%	4	&坐标旋转\\
%	5	&去均值与线性趋势	&去除零频与极低频率项\\
%	6	&尖灭			&由于FFT隐含周期化条件，需要平滑零化首位记录以连续衔接\\
%	7	&滤波\\	
%	\bottomrule
%	\end{tabular}
%\end{table}
以上操作中，“去仪器相应”、“去均值与线性趋势”、“尖灭”三项均从根本上改变了采样点的数值，其执行需要严格遵守一定的先后顺序。在确定先后顺序中，需要深刻理解时间--频谱关系及仪器工作原理，否则将产生不可预知的假象和错误。一般预处理中，按照“去均值与线性趋势”-->“尖灭”-->“去仪器相应”的先后顺序进行。而“时间截断”和“坐标旋转”可以在后续操作中任意进行。“滤波”则针对研究目的的需要而进行。将预处理完的数据应用到研究中，如果涉及到频谱操作或（反）卷积操作，则需要再做一次“去均值与线性趋势”与“尖灭”。\par
当然，以上的步骤与顺序介绍并非放置四海而皆准。微小的预处理差异会带来截然不同的结果，所以需要备份数据，反复操作对比，并针对研究目的确定最合适的预处理流程和参数选择。\par
SAC软件预处理命令范例如下：\par

\begin{SACCode}
SAC> r eg.sac           #读入SAC文件
SAC> rmean              #去均值
SAC> rtr                #去线性趋势
SAC> taper ***          #尖灭
SAC> transfer           #去仪器响应
SAC> rotate to gcp      #旋转至大圆路径，即RTZ分量
SAC> cut t1 -50 200     #截取t1时间前50s到后200s
SAC> bandpass ***       #带通滤波
SAC> lowpass ***        #低通滤波
SAC> highpass ***       #高通滤波
SAC> bandrej ***        #带阻滤波
\end{SACCode}

%\clearpage\section{总结及Seismic-Field-Works操作流程简介}

\end{CJK}
\end{document}
